%% !TEX root = ../manual_gaolv5.tex
% Chapter 4.
%% Chapter 4: The rounding direction and the floating-point environment.
%% Sources: gaol.tex of f5bf180 lines 1109-1749 (the rounding direction, its
%% effects on the program, the NEON intervals, flush-to-zero, the
%% floating-point exceptions), 2301-2408 (initialization and cleanup),
%% 6124-6155 (the rounding functions of the FPU) and 475-598 (the AVX-512
%% path, in the options of CMake); doc/using.md lines 228-290 (initialization
%% and cleanup) and 1070-1972 (the rounding direction, the AVX-512 path, the
%% NEON and the FPU intervals, the control registers of 32-bit x86,
%% flush-to-zero, the floating-point exceptions).

\chapter{The rounding direction and the floating-point environment}
\label{chap:rounding}

An interval library computes its bounds with the rounding direction of the
processor set outward, and that direction is a state of the processor which
the whole program shares. How the library takes charge of this state, and
what it leaves to the program, decides both the correctness of the bounds and
the behaviour of every double the program computes beside its intervals. GAOL
4 left the direction to the user. GAOL v5 sets it itself, in each operation,
and the choice has consequences the user needs to know: the doubles of the
program are rounded upward while GAOL is initialized, the C library converts
its numbers in that direction, and the compiler restores nothing.

\medskip
The rounding direction is not the only state of the floating-point
environment that matters to the bounds. The modes that flush the subnormal
numbers to zero, which a program linked with \option{-Ofast} gets without
asking, make the bounds of some operations wrong, and the floating-point
exceptions, when a program enables them, stop the program in the middle of an
operation of GAOL. The chapter describes the three, and how GAOL defends its
bounds against them.

\medskip
In this chapter, we describe the rounding direction as GAOL v5 handles it,
and what it does to the rest of the program. The first section states how
each operation sets the direction, and which operations leave it alone. The
second section shows what the upward direction does to the doubles of the
program, and what a program does about it; the third describes the build that
preserves the direction of the program. The fourth, fifth and sixth sections
describe the AVX-512 path, which carries its rounding direction in its
instructions, and the NEON and the FPU intervals; the seventh describes the
control registers of the 32-bit x86 processors. The following two sections
treat the flush-to-zero modes and the
floating-point exceptions. The chapter ends with the reference of the
initialization and the cleanup of GAOL, and of its rounding functions. The
design of the kernels, which instructions compute the bounds and what the
check of the direction costs on each processor, is the subject of
Chapter~\ref{chap:computing} (\pxref{sec:kernels}); the options of the builds
are those of Chapter~\ref{chap:installation} (\pxref{sec:cmake-options}).

%%------------------------------------------------------------------------
\section{The rounding direction}\label{sec:rounding-direction}
%%------------------------------------------------------------------------
\cindex{rounding!direction}%
\cindex{rounding!downward}\cindex{rounding!upward}%
\cindex{F@\FSet}\cindex{R@\RSet}%
\cindex{containment property}%
\cindex{rounding!outward|see{outward rounding}}%
\cindex{downward rounding ($\roundDn{\:}$)}%
\cindex{upward rounding ($\roundUp{\:}$)}%
\cindex{outward rounding}%

Floating-point interval arithmetic requires \dfn{outward rounding} to fulfil
the containment property (\pxref{sec:containment}): the sum of \itv{a}{b} and
\itv{c}{d} is computed as \itv{\roundDn{a+c}}{\roundUp{b+d}}, where
\roundDn{r} is the greatest double of \FSet{} below the real $r$ and
\roundUp{r} the smallest one above it (\pxref{sec:rounding-definitions}). The
processor performs these two roundings when its rounding direction is set
towards $-\infty$ and towards $+\infty$ respectively.

\medskip
On most platforms, switching the rounding direction is costly. The number of
switches is cut down with the following property.

\begin{Prop}\label{prop:rounding-negation}
For every real $r$, $\roundDn{-r}=-\roundUp{r}$.
\end{Prop}

\noindent
Nearly every downward rounding is therefore replaced by an upward one, by
negating twice the operation performed. GAOL stores the lower bound of an
interval negated, and computes both bounds rounding upward: the rounding
direction is set towards $+\infty$ once, and no other direction is needed
(\pxref{sec:kernels}).

\medskip
GAOL~4 left the rounding direction to the user, who had to set it towards
$+\infty$ before any computation on intervals: this was its \dfn{trust
rounding mode}. GAOL v5 takes the direction in charge.\newinvfive{} Each
operation that computes a bound checks the direction, and sets it upward when
it is not, whatever the build; the operations of the AVX-512 path, which carry
their direction in the instruction, are the exception (\pxref{sec:avx512}).
The direction is then left upward, unless GAOL is built to preserve it
(\pxref{sec:preserve-rounding}). The bounds are therefore right whatever
rounding direction the calling code left.
\cindex{trust rounding mode}%

\medskip
The check is an addition, $1+2^{-60}$ or one of its variants below, which is above 1 only when rounded
upward, rather than a reading of the rounding direction, which costs far more
on some platforms: under Rosetta~2, reading it made an addition of intervals
take 800~ns rather than 2.8~ns, and it cost far more with 32-bit Visual C++
as well. The check takes three forms, by processor:
\begin{itemize}
\item on x86-64 processors, and on ARM processors with the FPU intervals and
  GCC or Clang, it is
  $1+(2^{-1060}+0)$, which shows the modes that flush the subnormal numbers to
  zero as well (\pxref{sec:ftz});
\item with the NEON intervals of 64-bit ARM processors (\pxref{sec:neon}), it
  is a read of the control register FPCR, which holds both the direction and
  these modes;
\item on 32-bit x86 processors, whose two units that round, the x87 unit and
  the SSE instructions, each have a direction, it is a read of the control
  word of the first and of MXCSR, with \code{fnstcw} and \code{stmxcsr}
  (\pxref{sec:x86-control-registers}).\newinvfive
\end{itemize}
\noindent
Each function of intervals checks the direction once, and the functions it
calls on doubles take it to be upward already. The cost of the check on each
processor is measured in Chapter~\ref{chap:computing}
(\pxref{sec:rounding-by-processor}).

\medskip
Unless GAOL is built to preserve it, the rounding direction stays upward from
the initialization of GAOL on, for the whole program, and \code{cleanup()}
sets back the one that the initialization found: every program calls
\code{gaol::cleanup()} right after its last use of GAOL
(\pxref{chap:init}).\newinvfive{} The direction set back is the one the first
\code{init()} found when GAOL initialized itself: to nearest, as a program
starts; on x86, for the x87 unit and for the SSE instructions each; in the
thread calling \code{cleanup()}. The rest of the floating-point environment is
left as it is, the exception flags raised in particular.
\cindex{rounding!preserving}%

\medskip
Some operations set no direction, and leave the one they found, with or
without \option{GAOL\_PRESERVE\_ROUNDING}: a program that wants a direction
after one of them sets it itself.\newinvfive
\cindex{rounding!left as found}%
\begin{itemize}
\item Those that round nothing and make no check: the negation and the unary
  plus, \code{abs()}, \code{min()}, \code{max()}, \code{sign()}, the roundings
  to an integer (\code{floor()}, \code{ceil()}, \code{trunc()},
  \code{integer()}, \code{round\_ties\_to\_even()}\ldots), \code{\&} and
  \code{|}, \code{hull()} and \code{intersect()}, \code{invabs\_rel()},
  \code{mig()} and \code{mag()}, \code{nth\_root(x, 0)} and
  \code{nth\_root(x, 1)}, the constructors from bounds, the access to the
  bounds, \code{is\_empty()}, the relations and the other predicates,
  \code{exact\_string()}.
\item Those that return before their check. Some return the empty set at once
  for an empty operand, and some do not: \code{y * empty} and
  \code{exp(empty)} leave the direction, where \code{empty + y} and
  \code{sqrt(empty)} set it upward, so that a program counts on neither. And
  \code{hausdorff(x, y)} returns $+\infty$ when a bound is infinite in one of
  the two intervals only, where it sets the direction upward otherwise.
\item The operations of the AVX-512 path, which set none
  (\pxref{sec:avx512}).
\end{itemize}

%%------------------------------------------------------------------------
\section{What the upward rounding does to the program}
\label{sec:upward-effects}
%%------------------------------------------------------------------------
\cindex{rounding!effects on the program}%

The rounding direction is the processor's, and the whole program shares it.
The results of GAOL do not depend on it, since each operation sets it upward
when it is not, in every thread; the doubles the program computes beside its
intervals do.\newinvfive{} Unless GAOL is built with
\option{GAOL\_PRESERVE\_ROUNDING}, they are rounded upward from the
initialization of GAOL, before \code{main()}, until \code{cleanup()}, and
again after each operation of GAOL that follows it and sets the direction. In
a build that takes the AVX-512 path, which without
\option{GAOL\_PRESERVE\_ROUNDING} only the developers ask for, the arithmetic
operations set none, and the direction after one of them is the program's.

\medskip
The compiler does not undo it. With the flags of \code{gaol::gaol},
\option{-frounding-math} in particular, an inexact operation on constants
written in a function, \code{1.0 / 3.0}, is not evaluated when compiling, to
nearest, but when the program runs, in the direction in effect. Only what the
compiler must evaluate when compiling, a \code{constexpr} or the constant
initializer of a static variable, is still rounded to nearest.

\medskip
Table~\ref{tab:upward-effects} gives what the program computes, with GCC~9.4
or Clang~18 and the GNU C library~2.31 on x86-64;
\file{examples/13\_rounding\_environment.cpp} of the sources computes each row
and checks it.

\begin{table}[htbp]\centering
\renewcommand{\arraystretch}{1.1}\footnotesize\setlength{\tabcolsep}{3pt}
  \caption[What the program computes while GAOL is initialized and after \code{cleanup()}]{What the program computes while GAOL is initialized and after \code{cleanup()}, with GCC~9.4 or Clang~18 and the GNU C library~2.31 on x86-64.} \label{tab:upward-effects}
  \begin{tabular}{L{.34\textwidth}L{.29\textwidth}L{.29\textwidth}}
    \hline
    The program computes & While GAOL is initialized & After \code{cleanup()}\\
    \hline\hline
    \code{1.0 / 3.0} & 0.33333333333333337 & 0.33333333333333331\\
    \code{std::stod("0.3") == 0.3} & false & true\\
    \code{std::lrint(2.3)} & 3 & 2\\
    \code{std::printf("\%.2f", 2.675)} & 2.68 & 2.67\\
    \code{std::cout <{}< std::setprecision(3) <{}< 3.14159265358979} & 3.15 & 3.14\\
    0.1 added $10^7$ times & 1000000.0005483569 & 999999.99983897537\\
    20\,000 doubles written with \code{\%.17g}, read back by \code{std::strtod} & 20\,000 come back different & none\\
    \code{TwoSum(1e20, 1e-20)}, the error of a sum & $s+e\neq a+b$ & $s+e=a+b$\\
    \code{TwoSum(a, b)}, 50\,000 pairs of doubles of exponents from $-60$ to 60 & $s+e\neq a+b$ for 10\,717 pairs & none\\
    \code{TwoProd(a, b)}, Dekker's splitting, 50\,000 pairs of doubles of $[1,2)$ & $p+e\neq ab$ for 1\,888 pairs & none\\
    \code{TwoProd(a, b)}, \code{e = std::fma(a, b, -a*b)}, the same pairs & none & none\\
    \hline
  \end{tabular}

\end{table}

The quotient is the double above $1/3$, where the program expects the
nearest. The C library rounds in the direction in effect too, as the C
standard recommends: with the GNU C library, \code{std::rint(-2.7)} is $-2$
where it is $-3$ to nearest, and \code{strtod()}, \code{stod()},
\code{printf()} and the streams round up the decimal digits they read and
write. A double written with 17 digits, which reads back to itself to nearest,
then comes back one double above, and one more at each round trip. Another C
library may convert to nearest whatever the direction: the example reports
what the one it runs with does. Sums are biased: each addition rounds up, so
0.1 added ten million times ends $5.5\cdot10^{-4}$ above the exact
$10^7\times 0.1$, where to nearest it ends $1.6\cdot10^{-4}$ below.

\medskip
The error-free transformations break. TwoSum, the algorithm that recovers the
error of a sum,

\begin{example}
s = a + b;  bb = s - a;  e = (a - (s - bb)) + (b - bb);
\end{example}

\noindent
gives $s+e=a+b$ only to nearest. Rounded upward, the exact error of
$10^{20}+10^{-20}$, which is $10^{-20}-16384$, is no double: \code{s} is
$10^{20}+16384$ and \code{e} is $-16383.999999999998$, whose sum is not
$10^{20}+10^{-20}$. TwoProd, which recovers the error of a product from
Dekker's splitting of the operands, breaks the same way: 1\,888 of the
50\,000 pairs of the table, drawn in $[1,2)$, lose their exactness upward,
where pairs of widely different magnitudes keep it, the error of the product of
the split parts they drop staying below the error of the rounded product. It
is exact in every direction when it is computed with \code{std::fma} instead:
the error of a rounded product is a double, which \code{std::fma} computes
however it rounds, provided it is correctly rounded, by the instruction or by
a call of a C library whose \code{fma()} is (with GCC, \code{std::fma} without
optimization is a call of the C library, whose \code{fma()} is wrong on
Cygwin, where \code{\_\_builtin\_fma} is the instruction with
\option{-mfma}, which GAOL calls itself). Compensated sums, and the
double-double arithmetic built on such algorithms, lose their accuracy the
same way. Built with \option{GAOL\_PRESERVE\_ROUNDING}, the program computes
the right-hand column of the table all the time.

\medskip
\noindent
The upward direction has other consequences:
\begin{itemize}
\item a thread created after the initialization of GAOL starts upward too
  (\code{std::thread} gives a new thread the floating-point environment of the
  thread constructing it), and \code{cleanup()} sets back the direction of the
  calling thread only;
\item the libraries the program calls run upward as well: the \code{exp},
  \code{log} and \code{pow} of the GNU C library~2.31 go from 0.51 ulp to 1.01
  ulp from the exact value, its \code{cbrt} from 2.8 to 3.8 ulp (20\,000
  random arguments each, against mpmath), and its \code{sin}, \code{cos},
  \code{tan} and \code{atan} give the same results;
\item GCC does not model the rounding direction, even with
  \option{-frounding-math} (GCC bug~34678,
  \url{https://gcc.gnu.org/bugzilla/show_bug.cgi?id=34678}), and reuses after
  \code{cleanup()} a double computed before it, still rounded upward, where
  the program computes it again: from \option{-O1}, GCC~9.4 does it for
  \code{a / b} written on both sides of the call, where Clang~18 computes the
  quotient again. A program cannot rely on either.
\end{itemize}

\medskip
\noindent
What a program does about it:
\begin{itemize}
\item[$\bullet$] \textbf{It calls \code{cleanup()} right after the last use of
  GAOL}, before it prints, converts or sums what it computed: its doubles are
  then computed to nearest. A program that goes back to GAOL afterwards calls
  \code{restore\_rounding()} (or \code{std::fesetround(FE\_TONEAREST)}) after
  each use, since an operation of GAOL sets the direction upward again and
  \code{cleanup()} sets it back once only (\pxref{chap:init}).
\item[$\bullet$] \textbf{It reads and writes numbers through GAOL}, not
  through the C library: \code{gaol::textToInterval("0.3")} encloses $3/10$ in
  every direction, where \code{std::stod("0.3")} does not, and
  \code{exact\_string()} and \code{\%a} write a number that reads back to the
  same bits in every direction. GAOL writes its intervals rounded outward
  whatever the direction.
\item[$\bullet$] \textbf{It computes in a block the doubles that need rounding
  to nearest} (an error-free transformation, \code{std::lrint()}, a decimal
  conversion, a call of a library written for that): it saves the direction
  with \code{std::fegetround()}, sets \code{FE\_TONEAREST}, and sets the saved
  one back when the block ends, with no operation of GAOL inside, which would
  set the direction upward again. The example of the sources writes such a
  guard, \code{nearest\_scope}.
\item[$\bullet$] \textbf{It does not count on the compiler} to compute a
  double where the source code does. One that has to be computed on this side
  of a change of direction is computed from operands read again through
  \code{gaol\_core::rnd\_keep()}, which writes its argument to volatile memory
  and returns it: \code{gaol\_core::rnd\_keep(a) / gaol\_core::rnd\_keep(b)}
  after \code{cleanup()} is the quotient to nearest with GCC too, where
  \code{gaol\_core::rnd\_keep(a / b)} may be the one computed before.
\item[$\bullet$] \textbf{Or GAOL is built with
  \option{GAOL\_PRESERVE\_ROUNDING}}, and the program keeps its own direction
  (\pxref{sec:preserve-rounding}).
\end{itemize}

%%------------------------------------------------------------------------
\section{Preserving the rounding direction}\label{sec:preserve-rounding}
%%------------------------------------------------------------------------
\cindex{rounding!preserving}%
\cindex{GAOL PRESERVE ROUNDING@\option{GAOL\_PRESERVE\_ROUNDING}}%

Code that needs its own rounding direction after the operations of GAOL
builds GAOL with the option \option{GAOL\_PRESERVE\_ROUNDING} of CMake, which
is \loption{enable-preserve-\allowbreak rounding} with autotools and
\option{-Denable-preserve-\allowbreak rounding=true} with meson
(\pxref{sec:cmake-options}). It is off by default in the three builds. With
it, each operation reads the rounding direction it finds, sets it upward, and
sets back the one it found when it is done; it also sets back the modes
flushing the subnormal numbers to zero that it cleared (\pxref{sec:ftz}). The
operations of the AVX-512 path set none (\pxref{sec:avx512}). Nothing else
changes: the bounds are the same, and the program still calls
\code{cleanup()}.\newinvfive

\medskip
\noindent
The option changes three other things.
\begin{itemize}
\item[$\bullet$] \textbf{The AVX-512 path.} With it, the build takes the
  AVX-512 path (it was the default of an option,
  \option{GAOL\_PREFER\_AVX512}): where the machine building has the AVX-512
  instructions, the four arithmetic operations, the square, the square root
  and, with FMA, the integer powers beyond 1 and $-1$ take them, which carry
  their rounding direction, set none and store none. The library then needs a
  processor with them.\newinvfive
\item[$\bullet$] \textbf{The initialization.} It leaves the whole
  floating-point environment as it found it, where the default build sets the
  default environment: an exception a static object enabled before it stays
  enabled (\pxref{sec:fp-exceptions}).
\item[$\bullet$] \textbf{\code{restore\_rounding()}.} It does nothing, GAOL
  leaving the direction of the program alone (\pxref{chap:init}).
\end{itemize}

\medskip
\noindent
The direction is written only where it has to be. On the processors whose
control registers GAOL writes itself, an operation reads the direction and the
flush modes once at its start, writes the register only if the direction is
not upward or a flush mode is set, and writes back what it found at its end
only if it differs from what the operation left: a program that computes
upward already, with the flush modes clear, pays no write. These registers
are FPCR on 64-bit ARM and FPSCR on 32-bit ARM with GCC (\pxref{sec:neon}),
and the control word of the x87 unit and MXCSR on x86
(\pxref{sec:x86-control-registers}).

\medskip
The option costs time, less with the AVX-512 path, which it turns on where
the machine has the instructions. Table~\ref{tab:preserve-cost} gives the
factor by which each operation gets longer on an Intel i7-1185G7 with
Clang~21.1.8 (\file{tests/performance.cpp} and the long lines of
\file{doc/compare/code/run\_lines.sh}, the median of 9 interleaved runs of
6 November 2026), with the path, and without it (a build for a machine without
the instructions, or one that refuses the path with
\option{-DGAOL\_PREFER\_AVX512=OFF}, a setting of the developers). The
arithmetic operations take 1.2 to 1.6 times longer with the path, and several
times longer without it; a line of 15 additions of intervals takes 2.1 times
longer with the path and 8.5 times without.

\begin{table}[htbp]\centering\footnotesize
  \caption[The cost of \option{GAOL\_PRESERVE\_ROUNDING}, with and without the AVX-512 path]{How many times longer each operation takes with \option{GAOL\_PRESERVE\_ROUNDING} than without it, with and without the AVX-512 path (Intel i7-1185G7, Clang~21.1.8, \file{tests/performance.cpp}, the median of 9 interleaved runs, 6 November 2026).} \label{tab:preserve-cost}
  \begin{tabular}{lcc}
    \hline
    Operation & With the AVX-512 path & Without it\\
    \hline\hline
    \code{x + y}, \code{x - y} & 1.2 & 7.2\\
    \code{x * y} & 1.6 & 4.3\\
    \code{x / y} & 1.6 & 3.4\\
    \code{sqr(x)} & 1.0 & 4.4\\
    \code{sqrt(x)} & 0.6 (it is faster) & 1.5\\
    \code{pow(x, 3)} & 1.2 & 1.6\\
    \code{exp(x)} & 1.2 & 1.2\\
    \code{log(x)} & 1.1 & 1.0\\
    \code{sin(x)}, \code{cos(x)} & 1.2 & 1.1 to 1.2\\
    a line of 15 additions & 2.1 & 8.5\\
    \hline
  \end{tabular}

\end{table}

%%------------------------------------------------------------------------
\section{The AVX-512 path}\label{sec:avx512}
%%------------------------------------------------------------------------
\cindex{AVX-512}%

The AVX-512 instructions of the x86-64 processors carry a rounding direction
in the instruction itself, the \dfn{embedded rounding}: an operation that
takes them reads no direction, sets none and stores none. GAOL takes them in
the build that preserves the rounding direction, in which every operation of
the other path saves and restores the direction
(\pxref{sec:preserve-rounding}): the path relieves \code{+}, \code{-},
\code{*}, \code{/}, \code{sqrt}, \code{sqr} and the integer powers of that
store.\newinvfive{} The instructions of the kernel are described in
Chapter~\ref{chap:computing} (\pxref{sec:kernels}).

\medskip
\noindent
\textbf{When the build takes it.} The path is decided when GAOL is configured,
as the fused multiply-add instructions are (\pxref{sec:cmake-options}): where
\option{GAOL\_PRESERVE\_ROUNDING} is on, the compiler takes the instructions
and the machine building runs them, the macro \code{GAOL\_AVX512\_ARITHMETIC}
is defined in \file{gaol/gaol\_configuration.h}; a cross-compilation leaves
the path out, and the operations test nothing as they run. With that option
every operation gains, 2.8 times at the median of seven operations, from 1.3
to 4.4 times, on an Intel i7-1185G7, whose AVX-512 units are halved; without
it they cost, 17\,\% at the median, the direction being written only when it
is not upward already, so the path is left out. \option{GAOL\_PREFER\_AVX512},
\loption{enable-prefer-avx512} and \option{-Denable-prefer-avx512} were
options of the user, and are now settings of the developers:\newinvfive{}
\code{=ON} asks for the path without the preserved direction, and
\option{-DGAOL\_PREFER\_AVX512=OFF}, \loption{disable-prefer-avx512} and
\option{-Denable-prefer-avx512=false} refuse it, for a package for any
processor of the architecture. A library built with the path needs a processor
with the instructions, as with the fused multiply-add instructions, and
\code{gaol::init()} stops a program run on another one, with a message on the
standard error, rather than letting its first operation stop on an illegal
instruction (an operation made before GAOL initializes itself, as described
in \exref{chap:init}, still does). GAOL no longer chooses the path as the
program runs, in \code{init()}, with every operation testing the flag it set,
in the default build too.\newinvfive

\medskip
\noindent
\textbf{The operations.} On such a build, \code{x + y}, \code{x - y},
\code{x * y}, \code{x / y}, \code{sqrt(x)} and \code{sqr(x)}, and the
operations \code{+}, \code{-}, \code{*} and \code{/} with a double, are
computed by the rounding direction the instructions carry. So are
\code{pown(x, n)} and \code{pow(x, n)} for an integer $n$ of 2 or more
($n=2$ is \code{sqr(x)}), or of $-2$ or less where \code{x} has no zero,
whose two corners are within the range of the exact products of their kernel;
the exponents 0, 1 and $-1$, and the powers beyond the range, which a corner
near 0 or a large exponent makes, take the path of the direction. These
operations set no rounding direction: the program keeps its own, in the
default build as with \option{GAOL\_PRESERVE\_ROUNDING}, and the other
operations (\code{\%}, the elementary functions, the reading of a number, the
integer powers beyond the range) leave it upward, as they always do.

\medskip
\noindent
\textbf{The flags and the flush modes.} The instructions of the path raise no
floating-point exception flag, the architecture requiring the suppression of
the exceptions with the embedded rounding; the check of the other operations
raises the flags described in \exref{sec:fp-exceptions}. The modes that flush
the subnormal numbers to zero, which the embedded rounding honours as the
other instructions do (measured on an Intel i7-1185G7: with flush-to-zero
set, the difference of \code{3e-308} and \code{2.9e-308} is 0), are cleared
before the operation reads a bound when a program or a plug-in has set them,
as the other operations clear them (\pxref{sec:ftz}), and set back after its
result with \option{GAOL\_PRESERVE\_ROUNDING}: with Clang~18, the difference
of \code{x - y} was once computed after flush-to-zero had been set back, and
\code{[3e-308] - [2.9e-308]} was \itv{0}{0}.\newinvfive

\medskip
\noindent
\textbf{The bounds.} They are those of the other path, bit for bit, but for
the sign of a zero bound in one case.

\begin{Rem}\label{rem:avx512-zero}
\code{x * d}, \code{x *= d} and \code{d * x}, which is computed as
\code{x * d},\newinvfive{} of an interval whose zero bounds have different
signs, \itv{+0}{-0}, are \itv{-0}{+0} with the AVX-512 path and \itv{+0}{-0}
with the SSE2 one, $+0$ and $-0$ being the same number
(\pxref{sec:zero-sign}). The rounding rule is the same, and the tests of
\file{tests/}, which a job of the continuous integration runs with the option
on, check it with \code{==}, which does not see the sign of a zero;
\code{rounding\_direction} checks the bits of that one case in each path, and
every result of the operations of intervals and with a double on the
intervals of two bounds among 0 of either sign, $\pm1$, $\pm2.5$, $\pm3$,
$\pm10^{-310}$, $\pm10^{300}$ and $\pm\infty$ was compared bit for bit.
\end{Rem}

\medskip
\noindent
\textbf{The cost.} The operations of the path cost more time on a processor
that halves its AVX-512 instructions, where they cost less on one that does
not. On an Intel i7-1185G7 (halved), with Clang~21.1.8
(\cmd{gaol\_performance}, the median of 9 runs of 26 October 2026, from
Table~\ref{tab:computing-avx512}), \code{x * y} takes
3.7~ns rather than 2.3 in the default build, \code{x / y} 4.1~ns rather than
2.6 and \code{x + y} 1.5 rather than 1.1, but \code{sqrt(x)} 4.4 rather than
8.8, the instruction rounding each root where the SSE2 path checks the root
of the C library with a product. With \option{GAOL\_PRESERVE\_ROUNDING}, where
the SSE2 path saves and restores the direction, \code{x + y} takes 1.9~ns
rather than 8.4, \code{sqr(x)} 2.0 rather than 8.3 and \code{pown(x, 3)} 7.5
rather than 9.9; the comparison of the libraries (\pxref{sec:comparison}) has
the preserved column of the other operations. The processors that do not
halve them, as the Xeon of the continuous integration, are expected to take
less. Which of the paths a GAOL took,
\file{examples/13\_rounding\_environment.cpp} tells.

%%------------------------------------------------------------------------
\section{The NEON intervals}\label{sec:neon}
%%------------------------------------------------------------------------
\cindex{NEON}\cindex{FPCR}\cindex{FPSCR}%

On 64-bit ARM processors (AArch64), GAOL computes its intervals with NEON
instructions, as it does with SSE2 instructions on x86.\newinvfive{} The two
bounds of an interval are the two lanes of one \code{float64x2\_t}, the lower
bound negated as in the \code{\_\_m128d} of the SSE2 intervals, and \code{+},
\code{-}, \code{*}, \code{/}, \code{sqr()} and \code{sqrt()} are kernels that
take and give their bounds in registers. The sum and the difference are one
vector addition rounded upward; the product is the maximum of the four
products of the bounds, with no branch on their signs; the quotient is one
vector division when the divisor has no zero, and single divisions otherwise;
the square roots of both bounds are one vector square root rounded upward,
the lower bound being that root where it is exact, and the double below it
otherwise (\pxref{sec:kernels}).

\medskip
They are the intervals of GAOL on 64-bit ARM with every compiler (GCC, Clang
and AppleClang under Linux and macOS, clang-cl and the Clang of MinGW-w64 and
MSYS2, and Visual C++), by default: the option is \option{GAOL\_SIMD} of
CMake, \loption{enable-simd} of \cmd{configure} and \code{enable-simd} of
meson, as for the SSE2 intervals (\pxref{sec:cmake-options}), and the macro
\code{GAOL\_USING\_NEON\_INSTRUCTIONS} is then 1 in
\file{gaol/gaol\_configuration.h}. They need no flag beyond those of interval
arithmetic (\pxref{sec:compiling}). \option{GAOL\_ARM64},
\loption{enable-arm64} and \option{-Denable-arm64} were an option of their
own, independent of \option{GAOL\_SIMD}, and are gone.\newinvfive{} 32-bit
ARM, whose NEON has no vector of doubles, does not take them: the option is
ignored there, and the FPU computes the intervals, as before
(\pxref{sec:fpu-intervals}). A universal build of macOS made by CMake takes
them for its arm64 part, and the SSE2 intervals for its x86\_64 part.

\medskip
Visual C++ writes the control register FPCR through the intrinsics
\code{\_ReadStatusReg()} and \code{\_WriteStatusReg()}, where GCC and Clang
have \code{mrs} and \code{msr} in an asm statement; it has no asm statement
to order its operations with respect to those writes, and relies on
\option{/fp:strict} to keep them where the source puts them, which
\file{tests/rounding\_direction.cpp} checks. With \option{/W4}, Visual C++
for arm64 warns (C4324, ``structure was padded due to alignment specifier'')
of each structure that holds an interval, its \code{float64x2\_t} being
aligned on 16 bytes there, and of \code{expr\_eval} of the installed
\file{gaol/parser/gaol\_expr\_eval.h}: the warning is harmless, and no job of
the continuous integration compiles with \option{/W4 /WX} for
arm64.\newinvfive

\medskip
\noindent
\textbf{The control register FPCR.} 64-bit ARM has no instruction that
carries its own rounding direction, as the AVX-512 path has: the direction is
the field RMode of the control register FPCR (bits 23 and 22, upward being
\code{0x400000}), which the scalar and the vector instructions share, and
which holds the modes that flush the subnormal numbers to zero as well (FZ,
bit 24, and from Armv8.7 FIZ, bit 0). The NEON intervals therefore check the
direction by reading FPCR with \code{mrs}, which raises no exception flag, and
write it with \code{msr} only if it is not upward with the flush modes clear,
where the FPU intervals add $1+(2^{-1060}+0)$, which raises the inexact flag
(\pxref{sec:fp-exceptions}). \code{msr} on FPCR is a serializing instruction,
slower than \code{mrs}, which is why it is written when needed only. With
\option{GAOL\_PRESERVE\_ROUNDING}, each operation reads FPCR once, and writes
back the direction and the modes it found at its end only if they differ from
what the operation left: a program that computes upward already, with the
flush modes clear, pays no write. \code{round\_upward()},
\code{round\_nearest()} and \code{round\_downward()} write FPCR as well, and
not through \code{fesetround()}, and only if the direction changes
(\pxref{chap:fpu}). \code{fegetround()} and \code{fesetround()} of the C
library read and write the same bits (the \code{FE\_*} values of the C
libraries of 64-bit ARM are these bits, which a static assertion checks), so
that a program may use either.

\medskip
\noindent
\textbf{The register FPSCR of 32-bit ARM.} The FPU intervals of 32-bit ARM
processors with GCC and a floating-point unit of doubles do the same for
\code{round\_upward()}, \code{round\_nearest()} and \code{round\_downward()},
and for the direction that \option{GAOL\_PRESERVE\_ROUNDING} restores
(\code{GAOL\_RND\_FPSCR\_DIRECT} of \file{gaol/gaol\_config.h}, 1 there;
\option{-DGAOL\_RND\_FPSCR\_DIRECT=0} gives the C library back). The register
is FPSCR, read with \code{vmrs} and written with \code{vmsr}, with the same
field RMode (the \code{FE\_*} values of the C libraries of 32-bit ARM are its
bits, which the same assertion checks). The flags of the exceptions and of
the comparisons are in it too, so that the functions that write it from the
value they read at once (\code{round\_nearest()}, \code{round\_upward()}, the
end of an operation) read it again just before, and no flag raised since is
lost. With \option{GAOL\_PRESERVE\_ROUNDING}, each operation reads FPSCR once
at its start, writes it only if the direction is not upward or a flush mode
is set, and writes back what it found at its end only if it differs. On the
Debian armhf jobs of the continuous integration, which run on the arm64
runners of GitHub, \code{x + y} took 22.5~ns rather than 49 with
\option{GAOL\_PRESERVE\_ROUNDING}, a line of 15 additions 318~ns rather than
721, and \code{midpoint()} 38~ns rather than 63.\newinvfive{} The check of
every operation without the option stays the addition: \code{vmrs} took
longer than the two additions there, \code{x + y} 6.1~ns rather than 4.1 and
a line of 15 additions 87~ns rather than 57, and the 32-bit ARM intervals
keep their time.

\medskip
\noindent
\textbf{The bounds.} They are those of the FPU intervals but for the sign of
a zero bound, which the maximum of 64-bit ARM, ranking $+0$ above $-0$, can
set otherwise than the SSE2 and the FPU intervals do, where the \code{maxpd}
of x86 returns its second operand (\pxref{sec:zero-sign}).

\begin{Rem}\label{rem:neon-zero}
The lower bound of \code{[-1, 0] * [-1, 0]} is $-0$ with the NEON intervals
and $+0$ with the SSE2 ones (the FPU intervals leave it to the compiler,
\pxref{sec:fpu-intervals}), and \code{interval::zero()} is stored
\code{<+0, +0>}, as with SSE2, so that its lower bound is $-0$; $+0$ and
$-0$ being the same number, \code{==} does not see the difference. The tool
\file{tests/tools/arith} compares the NEON and the FPU intervals on a grid of
operands, with the calling code in each of the four rounding directions: it
found no other difference. The same kind of caveat exists for the AVX-512
path (Remark~\ref{rem:avx512-zero}).
\end{Rem}

%%------------------------------------------------------------------------
\section{The FPU intervals}\label{sec:fpu-intervals}
%%------------------------------------------------------------------------
\cindex{FPU intervals}%

The intervals of the processors that have neither SSE2 nor NEON keep their
two bounds as two doubles, the lower one negated as in the SSE2 and the NEON
intervals, and compute with the scalar floating-point unit. They are the
intervals of 32-bit ARM, POWER, s390x, RISC-V and the other processors; of
Windows on 32-bit x86 with a compiler that takes no flag asking for the
aligned memory the SSE2 intervals need (MinGW-w64 with GCC and Clang takes
the SSE2 intervals there, with \option{-faligned-new=8}, and so do Visual C++
and clang-cl, whose intervals need no aligned memory, and 64-bit Windows);
and of the builds that ask for them (\option{-DGAOL\_SIMD=OFF},
\loption{disable-simd} or \option{-Denable-simd=false}), which leave
\code{GAOL\_USING\_SSE2\_INSTRUCTIONS} and
\code{GAOL\_USING\_NEON\_INSTRUCTIONS} at 0 in
\file{gaol/gaol\_configuration.h}.

\medskip
\code{x * y}, \code{x / y} and \code{x \% y} of a divisor without zero, and
\code{sqr(x)} take no branch on the signs of the bounds in the
source:\newinvfive{} GAOL~4 took one of nine cases of the signs for a
product, and went down a chain of up to eight tests for a quotient, which a
processor mispredicts when the signs come at random. The product is the
largest of four products for the upper bound and of four others for the lower
one, each with a factor negated, a product and its opposite not being rounded
the same way where the only direction is upward; the quotient divides each
bound of the dividend by the bound of the divisor that its sign chooses, two
divisions whatever the signs; the square is the largest of two squares for
each bound, the lower one being 0 for an interval that straddles 0. They are
the kernels of the SSE2 and the NEON intervals on two doubles each
(\pxref{sec:kernels}), and give the bounds of the SSE2 ones bit for bit for
the product of finite bounds and the quotient by a divisor without zero
(checked against a model of their kernels, not by running them), but for the
sign of a zero bound of a product (Remark~\ref{rem:fpu-zero}). An interval
with an infinite bound, whose 0 times an infinity is no product to form,
takes a function that sets both factors of such a product to 0 (the nine
cases of GAOL~4 took 12.4~ns for operands with an infinite bound drawn at
random, and it takes 9.2), and a divisor with a zero takes the table of the
NEON intervals, one division at most. \code{x - y} reads the bounds of
\code{y} after the check of the rounding direction, as \code{x + y} does: it
took three times as long as the sum, the bounds being kept across the calls
of the rare path of the check.

\medskip
\noindent
\textbf{The time of the kernels.} Table~\ref{tab:fpu-m2} gives the times
measured with \cmd{gaol\_performance} on an Apple M2 with AppleClang~21
(\option{-DGAOL\_SIMD=OFF}), the median of nine runs of each program, in turn
with the others, on 6 November 2026, in nanoseconds, with the code of GAOL~4
(the tree of 22 October 2026, before the kernels) and with the kernels of GAOL
v5.

\begin{table}[htbp]\centering\footnotesize
  \caption[The FPU intervals on an Apple M2 (\option{GAOL\_SIMD=OFF})]{The FPU intervals on an Apple M2 with AppleClang~21 (\option{-DGAOL\_SIMD=OFF}): the code of GAOL~4 and the kernels of GAOL v5, in nanoseconds (\cmd{gaol\_performance}, the median of nine runs).} \label{tab:fpu-m2}
  \begin{tabular}{lrr}
    \hline
    & GAOL~4 code & GAOL v5 kernels\\
    \hline\hline
    \code{x - y} & 3.54 & 1.15\\
    \code{x * y} & 3.25 & 3.25\\
    \code{x * y}, random signs & 6.39 & 3.31\\
    \code{x / y} & 3.98 & 3.23\\
    \code{x / y}, random signs & 3.17 & 3.21\\
    \code{sqr(x)}, random signs & 2.64 & 1.78\\
    a line of 7 \code{*} & 23.3 & 35.0\\
    a line of 7 \code{/} & 35.1 & 32.9\\
    a line of 13 \code{+ - * /} & 49.6 & 50.4\\
    \hline
  \end{tabular}

\end{table}

\noindent
The operands of \cmd{gaol\_performance} are 1\,024 intervals that the
processor learns by heart, so that a branch on their signs costs less than it
does on other operands. On 100\,000 intervals centered at random (the
benchmark of \file{doc/compare}, \pxref{sec:comparison}), \code{x * y} took
13.9~ns and takes 3.5, \code{x / y} 8.5 and 3.2, \code{sqr(x)} 6.9 and 1.8,
\code{x - y} 3.7 and 2.0, and Shekel 5\,298~ns and 138.

\medskip
A line of products of operands of one sign, which a processor predicts, takes
more time: 7 products in a line took 23.3~ns and take 35.0 (23.7 and 35.2 in
the benchmark of \file{doc/compare}, \code{pos\_mul8}; Apple M2). Each
product waits for the one before it, and the largest of four products waits
for its products and for two comparisons, where one product on a predicted
branch waits for one. The maximum of 64-bit ARM (\code{fmax}, one
instruction, which ranks $+0$ above $-0$) would shorten the wait: it is not
used, being a call with GCC for 32-bit ARM and x86. With
\option{GAOL\_PRESERVE\_ROUNDING} the write of the control register waits for
the result, so that the longer wait of the product shows there as well:
\code{x * y} takes 42.4~ns rather than 39.7 and the line of 7 \code{*} 297
rather than 279, that is 6 to 7\,\% more (\code{sqr(x)} 42.4 rather than
39.5; \code{x / y} 43.6 rather than 43.1; \code{x * y} of random signs 42.4
rather than 47.5, where the branches of the former code were mispredicted;
the median of nine runs of each code, in turn, on 6 November 2026). The operators write their result into
the bounds of \code{*this}, memory of the caller that is also given by address
to the functions of the rare cases, which a call or an asm statement with a
``memory'' clobber that restores the direction finds written: the products
they hold are computed before it, with the direction upward. That is how the
compilers treat it (GCC~15 for 32-bit ARM: the stores precede
\code{bl fesetround}), not a rule of the language, and no test holds it:
\code{GAOL\_RND\_KEEP()} would make it one, for 2.6~ns more with
\option{GAOL\_PRESERVE\_ROUNDING}, and the former product and \code{x + y}
did without.

\medskip
\noindent
\textbf{The processors of the continuous integration.} The processors and
compilers of the continuous integration that took the FPU intervals were
measured in turn (Visual C++ and clang-cl for arm64 take the NEON intervals
since),\newinvfive{} in one run of GitHub Actions, on an experiment branch:
each job builds the code before the kernels and after them, and runs
\cmd{gaol\_performance} five times with each, alternately, with the tools of
\file{.github/scripts/performance-ab.sh} and \file{compare-performance.py}
(\pxref{sec:scripts}). The compilers are GCC~13.3 and Clang~18.1 on
Ubuntu~24.04 for x86-64 (\option{-DGAOL\_SIMD=OFF}), GCC~14.2 on Debian~13
for 32-bit ARM, run natively on an arm64 runner, Visual C++~19.51 for x86 and
arm64, clang-cl~22.1 for x86 and arm64, and AppleClang~21 on macOS~26
(\option{-DGAOL\_SIMD=OFF}). Table~\ref{tab:fpu-ci} gives the ratio of the
median before the kernels to the median after them, so that a ratio above 1
means that the kernels are faster. These are figures of shared runners, of
one run, of the FPU intervals as they were: Visual C++ and clang-cl for x86
have the SSE2 intervals since,\newinvfive{} with their two bounds as two
doubles, and their rows are no longer those of the default build.

\begin{table}[htbp]\centering\scriptsize\setlength{\tabcolsep}{2.5pt}
  \caption[The FPU intervals on the processors of the continuous integration]{The FPU intervals on the processors of the continuous integration: the ratio of the median time of the code of GAOL~4 to the one of the kernels of GAOL v5 (five runs of \cmd{gaol\_performance} each, alternately; a ratio above 1 means that the kernels are faster). The columns are \code{x * y}, \code{x * y} of random signs, \code{x / y}, \code{sqr(x)}, \code{x - y}, a line of 7 \code{*}, a line of 7 \code{/} and a line of 13 \code{+ - * /}; ``preserved'' is the build with \option{GAOL\_PRESERVE\_ROUNDING}. $\dagger$: the spread of the five runs exceeds 1.15, and the ratio is noise; in bold, the ratios below 0.9 of a stable row.} \label{tab:fpu-ci}
  \begin{tabular}{l*{8}{r}}
    \hline
    & \code{x*y} & \code{x*y} rnd & \code{x/y} & \code{sqr} & \code{x-y}
      & 7 \code{*} & 7 \code{/} & 13 ops\\
    \hline\hline
    Ubuntu, GCC 13 & 1.29 & 1.64 & 1.25 & 1.09 & 1.02 & \textbf{0.58} & 1.13 & \textbf{0.76}\\
    Ubuntu, Clang 18 & 1.18 & 1.67 & 1.47 & \textbf{0.85} & 0.99 & 0.98 & 1.05 & 1.01\\
    armhf, GCC 14 & 1.41 & 2.28 & 1.63 & 1.17 & \textbf{0.86} & 1.11 & 1.25 & 1.04\\
    Visual C++, x86 & 1.68$^\dagger$ & 3.47 & 1.11 & 1.39 & 1.06 & 1.43 & 1.16 & 1.22$^\dagger$\\
    Visual C++, x86, preserved & 1.18 & 1.64$^\dagger$ & 1.02$^\dagger$ & 1.09$^\dagger$ & 1.02$^\dagger$ & 1.18 & 1.04 & \textbf{0.53}\\
    Visual C++, arm64 & 1.53 & 2.41 & \textbf{0.82} & 2.36 & 1.08$^\dagger$ & 1.71 & 1.05 & 1.41\\
    clang-cl, x86 & 1.06 & 1.23 & 1.24 & 1.06 & 1.04 & 1.04 & 1.16 & 1.05\\
    clang-cl, arm64 & \textbf{0.77} & 1.75 & 1.29 & \textbf{0.84} & 1.01 & \textbf{0.57} & 1.24 & \textbf{0.86}\\
    macOS, FPU & 0.92$^\dagger$ & 1.98$^\dagger$ & 1.13$^\dagger$ & 1.03$^\dagger$ & 2.49$^\dagger$ & 0.60$^\dagger$ & 1.04$^\dagger$ & 0.83$^\dagger$\\
    \hline
  \end{tabular}

\end{table}

\noindent
The product wins where the signs are random (1.23 to 3.47 times) and, alone,
of positive operands on 32-bit ARM, Visual C++ and Ubuntu. It loses on the
lines of products of positive operands with GCC on Linux (0.58 for the line
of 7 \code{*}, that is 1.7 times slower, where the Apple M2 above gives 23.3
and 35.0~ns, 1.5 times) and with clang-cl for arm64 (0.57), and alone with
clang-cl for arm64 (0.77): there, the branch of the former code is well
predicted and the kernel of eight products adds latency. The macOS runner is
too noisy to say (spreads of 1.30 to 2.71). The quotient wins from 1.11 to
1.63 times on six configurations, is within the noise on two (the direction
preserved, 1.02, and macOS, 1.13) and loses with Visual C++ for arm64 (0.82,
spread 1.04). \code{sqr(x)} is mixed: 0.84 with clang-cl for arm64 and 0.85
with Clang on Ubuntu, 2.36 with Visual C++ for arm64. \code{x - y} does not
change where the rows are stable (0.86 to 1.06), 0.86 on 32-bit ARM where
\code{x + y}, whose code did not change, gives 0.92 in the same job; the 3.54
to 1.15~ns of the M2 does not show. With the direction preserved on Visual
C++ for x86, the line of 13 operations is 1.9 times slower (0.53, spread
1.08), where each operation wins or stays: the cause was not searched. POWER,
s390x, RISC-V and MinGW-w64 for arm64 are not measured.

\medskip
What a compiler makes of the choices of the kernels is its own, and was read
in the assembly on Compiler Explorer, not run. The maxima of the product are
conditional moves with GCC (13 to 16) for 32-bit ARM, arm64 and x86, with
Clang, and with Visual C++ for x86 and arm64, \code{xsmaxcdp} instructions
with GCC~13 for POWER9, and branches on the products with GCC for POWER8,
s390x and RISC-V, which have no conditional move of doubles. The choices of
the quotient and of the square are conditional moves with GCC for POWER and
with Clang and Visual C++ for arm64, and branches on a sign with GCC for
32-bit ARM at \option{-O2} and \option{-O3} (two in the quotient with GCC 13
to 15, four with GCC~16, none in the square; at \option{-Os}, two in the
square with GCC~14 and~15), arm64 (one in the quotient), s390x and RISC-V,
and with Visual C++ for x86.

\begin{Rem}\label{rem:fpu-zero}
The sign of a zero bound of a product is not specified for the FPU
intervals: where two products are zeros of opposite signs, $+0$ and $-0$, the
maximum of the two is a comparison and a selection, and the compiler decides
which of the two it returns. The \code{maxsd} of x86 returns its second
operand, and the compilers that keep the order written do the same (the lower
bound of \code{[-1, 0] * [-1, 0]} was $+0$ on the jobs of the continuous
integration with GCC for 32-bit ARM and for arm64, with Clang for arm64 and
with Visual C++ for x86), but GCC~13 for POWER9 compiles
\code{(a > b) ? a : b} into \code{xsmaxcdp}, whose result for a $+0$ and a
$-0$ is its first operand (the lower bound was $-0$ on the POWER9 jobs under
QEMU, and the assembly of Compiler Explorer shows six \code{xsmaxcdp} in the
product and no other maximum; not run on a processor). The lower bound of
\code{[-1, 0] * [-1, 0]} is therefore $+0$ with the SSE2 intervals, $-0$ with
the NEON ones, and $+0$ or $-0$ with the FPU ones, by the compiler and the
target; so is the upper bound of \code{[-1, 0] * [0, 1]}. $+0$ and $-0$ are
the same number: \code{==} and \code{<} do not see it, and a program does not
read it (\pxref{sec:zero-sign}). The sign of the zero bound of a quotient and
of a square does not depend on a tie, and is fixed: \code{[0, 1] / [-2, -1]}
has the upper bound $-0$ (the zero bound of the dividend is divided by a
number whose sign the bounds of the divisor choose), and
\code{sqr(interval(-1, 2))} has the lower bound $+0$, which a selection of a
constant sets.
\end{Rem}

\medskip
The other bounds and the flags are those of the former code, as
\file{tests/tools/\allowbreak arith} (the four rounding directions of the caller) and 1.2
million random operands of every kind, with and without the mode that flushes
the subnormals to zero, show. \code{x * y} computes every corner product, as
the SSE2 and the NEON intervals do, so that a corner that is not a bound of
the result and underflows raises the underflow exception
(\code{FE\_UNDERFLOW}), which the former code did not, for 135\,808 of the
1.2 million \code{x * y}, and a corner that overflows raises the overflow
exception (\code{FE\_OVERFLOW}) for 23\,672 more, whose result has an
infinite bound all the same; no operation raises the invalid-operation or the
division-by-zero exception it did not (\pxref{sec:fp-exceptions}).

%%------------------------------------------------------------------------
\section{The control registers of 32-bit x86 processors}
\label{sec:x86-control-registers}
%%------------------------------------------------------------------------
\cindex{MXCSR}\cindex{x87 control word}%

A 32-bit x86 processor has two units that round, each with a direction of its
own: the x87 unit, in the field RC of its control word (bits 11 and 10), and
the SSE instructions, which compute the doubles of GAOL, in the field RC of
MXCSR (bits 14 and 13), where the modes that flush the subnormals to zero are
too (FTZ, bit 15, and DAZ, bit 6). Both have to be upward: the square root
and the \code{fma()} of a C library, and a \code{long double}, follow the
direction of the x87 unit (the elementary functions of CORE-MATH read MXCSR,
as on x86-64, after having read the one of the x87 unit with
\code{fegetround()}: the one of glibc, and $-1$ with the C runtime of Windows
when the two differ), and the addition that shows the direction elsewhere
sees the SSE instructions only. Each operation therefore reads the two
registers, the control word with \code{fnstcw} and MXCSR with
\code{stmxcsr}, in one asm statement (inline assembly with Visual C++), where
it called \code{fegetround()} and read MXCSR.\newinvfive{} That call took
about 80~ns with Visual C++ for x86: \code{x + y} took 81 to 100~ns, and a
line of 15 additions 1\,211 to 1\,546~ns, in the step of the continuous
integration that runs \cmd{gaol\_performance}.

\medskip
The registers are written (\code{fldcw}, \code{ldmxcsr}) only when the
direction of one is not upward or a flush mode is set, and the one that
differs only; MXCSR is written once for the direction and the flush modes.
With \option{GAOL\_PRESERVE\_ROUNDING}, an operation reads them once at its
start, and writes back at its end what it found, in the register that differs
from what the operation left: a program that computes upward already, with
the flush modes clear, pays no write. The SSE2 operations (\code{+},
\code{-}, \code{*}, \code{/}, \code{sqr()}\ldots), which use no x87
instruction, keep MXCSR alone with \option{GAOL\_PRESERVE\_ROUNDING} and leave
the x87 unit as it is: they write it only if it is not upward with the flush
modes clear, and give back what they found only if it was
not.\newinvfive{} A program that computes upward already pays no write there
either, \code{x + y} taking 1.7~ns rather than 8.2 and \code{x * y} 2.5
rather than 8.3 (GCC~15.2, an Intel i7-1185G7, the best of 15 runs of a
loop; a program to nearest pays 0.2~ns more, 8.4 for \code{x + y}).
\code{round\_upward()}, \code{round\_nearest()} and \code{round\_downward()}
write the register whose direction is not the one they set (both, as
\code{fesetround()} does, when neither is), and \code{cleanup()} and
\code{restore\_rounding()} set back the direction that \code{init()} found in
each unit (\pxref{chap:init}). The values are those of the hardware
(\code{0x0800} for upward in the control word, \code{0x4000} in MXCSR), not
the \code{FE\_*} constants of \file{<fenv.h>}, which are \code{0x0400} and
\code{0x0800} for downward and upward with glibc and \code{0x0100} and
\code{0x0200} with the C runtime of Windows and with mingw-w64~13 and later.

\medskip
An x86-64 processor has the two units too, and GAOL reads and writes their
registers itself there, as above, where it restores a direction
(\option{GAOL\_PRESERVE\_ROUNDING}, \code{cleanup()},
\code{restore\_rounding()}, and the operations that set the direction to
nearest and back: \code{midpoint()}, \code{chi()}, the writing and the
reading of an interval), rather than call \code{fegetround()} and
\code{fesetround()},\newinvfive{} which write both registers whatever they
hold, and then MXCSR once more for the direction of the SSE instructions:
three writes and a call, where a program that computes to nearest needs one
write and a program that computes upward none. The register that differs is
written only, with \code{fldcw} or \code{ldmxcsr}, which keep the flags;
Visual C++ for x64, which has no inline assembly and uses no x87 unit, reads
and writes MXCSR alone. mingw-w64, the only one that still called them, gave
reasons of its own to stop: from version 13, its \code{fegetround()} joins the
fields RC of the two units by OR, which is the direction of both only where
they agree (to nearest and upward give upward,
\code{FE\_DOWNWARD | FE\_UPWARD} is \code{FE\_TOWARDZERO}), and its
\code{fesetround()} clears the six exception flags of MXCSR when it changes
its value, so that the flags of the program were lost, in the operations with
\option{GAOL\_PRESERVE\_ROUNDING}, and in \code{cleanup()} and
\code{restore\_rounding()} with it or without it.

%%------------------------------------------------------------------------
\section{Flush-to-zero and denormals-are-zero}\label{sec:ftz}
%%------------------------------------------------------------------------
\cindex{flush-to-zero}%
\cindex{denormals-are-zero}%
\cindex{subnormal numbers}%
\cindex{-Ofast@\option{-Ofast}}%
\cindex{crtfastmath.o@\file{crtfastmath.o}}%

\begin{Def}\label{def:ftz}
The \dfn{subnormal numbers} are the nonzero doubles below $2^{-1022}$ in magnitude.
Processors have modes that flush them to zero. On x86, the
\dfn{flush-to-zero} mode (FTZ) and the \dfn{denormals-are-zero} mode (DAZ) of
the SSE instructions, two bits of the control register MXCSR, make every
subnormal result a zero (FTZ) and read every subnormal operand as a zero
(DAZ), in the comparisons too. On ARM, the mode FZ of the control register
(FPCR, FPSCR) does both.
\end{Def}

\noindent
With one of them set, the bounds of every operation with a subnormal result
or operand miss the exact result, whatever the rounding direction:
the product \code{interval(1e-300) * interval(1e-20)}, the square
\code{sqr(interval(1e-160))} and the difference
\code{interval(3e-308) - interval(2.9e-308)} are all \itv{0}{0}.\newinvfive

\medskip
\noindent
\textbf{How a program gets them.} A program linked with one of the options
\option{-Ofast}, \option{-ffast-math} and \option{-funsafe-math-optimizations}
gets them: GCC and Clang then link \file{crtfastmath.o}, on x86 and on ARM
Linux, whose constructor sets them for the whole program. GCC links it last,
and it runs after GAOL initialized itself; Clang links it first, and the
initialization of a static GAOL, which follows, clears the modes again,
unless GAOL is built with \option{GAOL\_PRESERVE\_ROUNDING}, or is a shared
library, initialized before the program. The option \option{-fno-fast-math}
of \file{gaol.pc} and \code{gaol::gaol} does not prevent it: it keeps
\file{crtfastmath.o} out of the link when it comes after
\option{-ffast-math}, but not after \option{-Ofast} (GCC~9.4 and~13 and
Clang~18 link it all the same, and \option{-fno-fast-math} silences the
\code{\#error} of \file{gaol/gaol\_config.h} against \option{-ffast-math},
which \option{-Ofast} would raise), nor, with GCC~9.4 and~13, after the
option \option{-funsafe-math-optimizations}. Loading a plug-in or a Python
module built with \option{-Ofast} sets the modes too, when it is loaded,
where its compiler linked \file{crtfastmath.o} into it: Clang~18 does, GCC
did before version~13.

\begin{Rem}\label{rem:ofast}
Linking a program that uses GAOL with \option{-Ofast}, \option{-ffast-math}
or \option{-funsafe-math-optimizations}, or loading code built with them,
sets the modes, which makes the bounds wrong unless GAOL defends itself
against them, as below. A program that uses GAOL is linked without these
options, or with \option{-mno-daz-ftz}, and the code that needs them is
compiled apart from it.
\end{Rem}

\noindent
GAOL defends itself in two ways.
\begin{itemize}
\item[$\bullet$] \textbf{The check of each operation clears the modes.} On
  x86 processors, and on ARM processors with GCC and Clang and with the NEON
  intervals of Visual C++, each operation that computes bounds clears the
  modes when it starts, in the check that sets the rounding direction upward:
  the check is one comparison with a subnormal number, which shows the modes
  as well as the direction (\pxref{sec:rounding-direction}), or, with the
  NEON intervals, a read of FPCR, which holds both; the modes stay cleared
  for the rest of the program, as the rounding direction stays upward. The
  check comes before the operation reads, compares or copies a bound, the
  test of the empty set aside, and the compiler is kept from making a
  comparison before it, which it would do, not knowing that the check changes
  how a subnormal compares: with GCC~13 at \option{-O3}, \code{asinpi()}
  compared the bounds it had read with 0 before its check, and with the modes
  set \code{asinpi()} of $[100\cdot 2^{-1074}]$ was $[32\cdot 2^{-1074}]$,
  above the exact value. With \option{GAOL\_PRESERVE\_ROUNDING}, each
  operation sets the modes back as it found them, once it has made its
  result, its maxima and minima of bounds included. The bounds of every
  operation that makes the check are then the ones it gives with the modes
  cleared, the tightest ones, and the modes are cleared, or set back, after
  it;\newinvfive{} \file{tests/rounding\_direction.cpp} checks it for each of
  them. It is the only defence against a plug-in, which sets the modes at any
  time, and against the compilers and processors without
  \option{-mno-daz-ftz}.

  The check costs little. On an Intel Xeon of the Cascade Lake generation,
  with Clang~18, an addition with a subnormal operand and result took as long
  as one of normal numbers (1.2~ns in a chain of dependent additions), as on
  an Intel i7-1185G7, and the check as long as the one of the direction
  alone, $1+2^{-60}$, in front of the addition of two SSE2 intervals. Through
  the library (\cmd{gaol\_performance}, medians of 9 interleaved runs),
  \code{x * y} took 4.0~ns rather than 3.4~ns, \code{sqrt} 11.0~ns rather
  than 10.3~ns and \code{pow(x, 3)} 13.7~ns rather than 13.0~ns, and as long
  with the same check made of a normal number: the cost of its second
  addition, not of the subnormal; \code{x + y}, \code{x - y}, \code{x / y},
  \code{sqr}, \code{exp}, \code{log}, \code{sin} and \code{cos} stayed within
  the noise of the machine. Some x86 processors take a microcode assist, of
  the order of a hundred cycles, for an operation with a subnormal operand or
  result, which each operation would pay there. \cmd{gaol\_performance}
  (\file{tests/performance.cpp}) prints the time of the check and of the same
  check made of $2^{-60}$, a normal double: on the 68 jobs of the continuous
  integration that run on processors (x86-64, i386, arm64 and armhf, under
  Linux, Windows and macOS, Rosetta~2 included), the first took 0.35 to
  1.43~ns and the second 0.40 to 1.38~ns, the two differing by $-0.63$ to
  $+0.44$~ns, either way from one job to another: none of these processors
  took such an assist. Under qemu (ppc64le, riscv64, s390x), both took 25 to
  83~ns.
\item[$\bullet$] \textbf{The link option \option{-mno-daz-ftz}.}
  \file{gaol.pc} and \code{gaol::gaol} give \option{-mno-daz-ftz} to the
  link, where the compiler accepts it: GCC~13 and later on x86, and from 11.4
  and 12.4 in the series 11 and 12 (GCC~9.4, 12.3, GCC for ARM and Clang~18
  refuse it, and the builds test the compiler with a link, not with its
  version). It keeps \file{crtfastmath.o} out of the link, so that a program
  linked with \option{-Ofast} does not get the modes at all, its own code
  included.
\end{itemize}

\medskip
\noindent
On the other processors (POWER, s390x, RISC-V, MIPS\ldots), whose FPU
intervals use the generic \code{round\_upward()} of
\file{gaol/gaol\_fpu\_fenv.h} (a call to \code{fesetround(FE\_UPWARD)}, whose
return value is not read), neither defence exists:\newinvfive{} the mode NI
(non-IEEE) of the FPSCR of POWER and the bit FS of the FCSR of MIPS, modes of
the kind of FTZ and DAZ that may flush the subnormal numbers, are not
cleared, and a program that sets them may get bounds that miss the exact
result for subnormal operands and results. What they do depends on the
processor, whether the toolchains of these processors set them in a program
linked with \option{-Ofast} was not checked, and nothing of GAOL was tested
with them set.

\medskip
\noindent
\textbf{The functions that make no check.} The functions that only read,
compare or copy the bounds of an interval make no check, which made each of
them 3 to 4.5~ns slower in a micro-benchmark on an Intel i7-1185G7, two to
three times the constructor or \code{certainly\_le()}: they compare the
bounds as bounds, whatever the modes.\newinvfive{} A mode that reads a
subnormal operand as a zero (DAZ, FZ) can make two different doubles compare
equal, both zero or subnormal, but never puts two doubles in the wrong order:
the doubles are compared first, and only where they compare equal are they
compared as integers, by their bits, which no mode changes. These functions
are the constructor from two bounds, the relations (\code{set\_contains()},
\code{certainly\_le()}, \code{set\_eq()}, \code{is\_zero()},
\code{lexicographic\_less}\dots), the intersection \code{\&} and the hull
\code{|}, \code{max()}, \code{min()}, \code{abs()}, \code{sign()},
\code{floor()}, \code{ceil()}, \code{integer()} and the other roundings to an
integer, \code{invabs\_rel()}, \code{mig()}, \code{mag()},
\code{is\_canonical()}, \code{feven()}, \code{nb\_fp\_numbers()}, and
\code{inf()} and \code{sup()} of \code{gaol\_ieee1788}; the roundings to an
integer round a subnormal bound from its bits, which the \code{floor()} and
\code{ceil()} of musl would return as it is. They give the results they give
with the modes cleared, and leave the modes as they found them, writing
nothing to the control register, with \option{GAOL\_PRESERVE\_ROUNDING} or
without. On the Intel i7-1185G7, with Clang~18 (\cmd{gaol\_performance},
medians of 9 interleaved runs), this left the constructor as fast as before,
and made \code{x <= y}, \code{x.set\_contains(y)}, \code{x \&= y} and
\code{x | y} 0.5 to 0.8~ns slower, and \code{floor()} 1.7~ns. Before, under
denormals-are-zero, \code{max()} of
$\itv{3\cdot2^{-1074}}{100\cdot2^{-1074}}$ and $[200\cdot2^{-1074}]$ was the
first one, \code{abs(interval(-1e-309, -1e-310))} stayed negative,
\code{interval(}$2^{-1073}$\code{, }$2^{-1074}$\code{)} kept its bounds in
the wrong order, and the hull \code{|} of the same two intervals, inline in
the headers of GAOL, depended on how the compiler of the program arranged its
comparisons.

\medskip
\code{midpoint()}, \code{split()} and \code{chi()}, which compute, and the
writing of an interval (\code{operator<{}<}, \code{intervalToText()}) make
the check, which clears the modes, or sets them back after them with
\option{GAOL\_PRESERVE\_ROUNDING} (1.3~ns more for \code{midpoint()}): the
midpoint of $\itv{2\cdot2^{-1074}}{4\cdot2^{-1074}}$ was 0, outside it, and
a subnormal upper bound was written rounded to nearest, below the bound,
where the C library writes a subnormal 0 under denormals-are-zero with gdtoa
(FreeBSD, macOS) and with the Debug runtime of Visual C++
(\file{tests/rounding\_direction.cpp} checks all of them).

\medskip
\noindent
\textbf{What remains.} The reader, which makes no check of the modes, takes
the numbers apart by their bits, and compares the two numbers of
\code{<a, b>}, which have to be the same double, by theirs too,\newinvfive{}
so that under denormals-are-zero \code{<1e-310, 1e-309>} is refused, as
without it (\pxref{sec:input-format}). The unary minus makes no check either,
but it only exchanges the stored bounds, which the modes do not change, as
they do not change the test of the empty set, made before the check of every
operation (an operation with an empty operand may return before its check,
the modes left as it found them), nor the test by which \code{hausdorff()}
returns $+\infty$ before its check when a bound is infinite in one of its two
intervals only. The modes of other processors, and of ARM with the FPU
intervals of Visual C++, are neither checked nor cleared, though the
functions above compare the bounds right there too, and \code{operator<{}<}
writes a subnormal bound outward, as $2^{-1022}$ away from zero and 0 toward
it; GCC links \file{crtfastmath.o} for none of the other processors GAOL is
tested on.

%%------------------------------------------------------------------------
\section{The floating-point exceptions}\label{sec:fp-exceptions}
%%------------------------------------------------------------------------
\cindex{floating-point exceptions}\cindex{exceptions!floating-point}%
\cindex{flags!floating-point}%

A program starts with the floating-point exceptions masked, and GAOL leaves
them as they are: an operation that raises one only sets its flag while the
exception is masked.

\begin{Rem}\label{rem:fp-exceptions-enabled}
A program that enables the floating-point exceptions, with
\code{feenableexcept()} of the GNU C library, \code{\_controlfp\_s()} of
Visual C++ or \code{fesetenv()} with an environment that traps, disables
them while GAOL computes, and enables them again afterwards if it wants them: the
operations of GAOL raise some, as listed below, and the processor then stops
the program (with \code{SIGFPE} on Linux) instead of letting the operation
give its bounds.\newinvfive
\end{Rem}

\noindent
What the operations raise:
\begin{itemize}
\item \code{log(interval(0,~1))}, which is \itv{-\infty}{0}, raises the
  divide-by-zero exception; \code{interval(1e308)*10}, which is
  \itv{\text{DBL\_MAX}}{+\infty}, raises the overflow exception.
\item Reading a number above the largest double (\code{1e400}, \code{1.8e308},
  \code{0x1p1024}) raises no overflow exception. The reader used to give it to
  \code{strtod()}, which computes $+\infty$ and raises it, so that
  \code{1e400*0}, which is \itv{0}{0}, raised the flag that stands for an
  infinite bound, and a program that had enabled the exception died while it
  read. The interval read is \itv{\text{DBL\_MAX}}{+\infty}; an operation on
  it raises the flag where its own result overflows (\code{1e400*2},
  \code{exp(1e400)}), as for any other interval. Nor does reading a number
  below the least normal double (\code{1e-310}, \code{1e-400},
  \code{0x1p-1075}) raise the underflow exception, which \code{strtod()}
  raises when it rounds such a number to a subnormal or to $0$.\newinvfive
\item Almost every operation raises the inexact exception: most bounds are
  rounded results, and where GAOL checks the rounding direction with an
  addition, $1+2^{-60}$ or $1+(2^{-1060}+0)$, whose result is inexact
  (\pxref{sec:rounding-direction}), even an operation whose bounds are exact
  raises it. On x86, the addition of the subnormal $2^{-1060}$ also raises
  the denormal-operand flag of the SSE instructions, which is none of the five
  flags of IEEE~754 (not in \code{FE\_ALL\_EXCEPT}), and its exception stops
  the program where the program enabled it; and with flush-to-zero set, the
  sum $2^{-1060}+0$, flushed to zero, raises the underflow exception
  (\code{FE\_UNDERFLOW}) too, before the check clears the mode. With the NEON
  intervals the check is a read of FPCR (\pxref{sec:neon}), which raises no
  flag: \code{+}, \code{-}, \code{/} and \code{sqrt()} raise the inexact one
  only when a bound they compute is rounded, \code{*}, \code{sqr()} and
  \code{pow(x, 2)}, which is \code{sqr()}, also when a corner product that is
  not a bound of the result is (\code{[-a, 128] * [-a, 4]} for
  $a = 1 + 2^{-52}$ has the exact bounds $-128a$ and 512, and the corner
  $a^2$ is rounded).\newinvfive
\item An empty operand raises no invalid-operation exception, though the
  bounds of the empty interval are NaN: every operation, relation, function
  and output of \code{gaol::interval} and of \code{gaol\_ieee1788} takes it
  without raising that exception. \code{is\_empty()} and the relations
  compare the bounds with the quiet comparisons of \file{<cmath>}
  (\code{std::islessequal()}\dots), each one instruction, as \code{<=} is,
  with GCC and Clang on x86-64, and so does the first comparison of the
  constructor \code{interval(l,~r)}, which gives the empty set for NaN bounds.
  With Visual C++, where these are calls to the C library, which made
  \code{x~*~y} and \code{sqrt()} 34 to 61\,\% slower on x64, GAOL compiles
  them itself into the same instruction, \code{ucomisd}, on x64 and on x86
  with SSE2; on arm64 they stay calls. With the invalid-operation exception
  enabled, \code{is\_empty()} of an empty interval killed the program, and so
  did about 50 operations with an empty operand, which compared its NaN
  bounds: \code{x~\&~y} for an empty \code{y}, \code{sqrt}, \code{exp},
  \code{min}, \code{max}, \code{floor}, \code{set\_contains()},
  \code{set\_disjoint()}, the output of the empty set\dots; and
  \code{interval::emptyset()} itself in a build without optimization. A NaN
  the program gives GAOL, the empty set for \code{interval(d)}, raises it no
  longer either: \code{interval(NAN)}, \code{interval(NAN,~1)},
  \code{x~=~NAN}, \code{x~+=~NAN} (and \code{-=}, \code{*=}, \code{/=},
  \code{\%=}), \code{x~\&=~NAN}, \code{x~<=~NAN} and \code{pow(x,~NAN)}
  compared it with \code{<} or \code{<=}. An empty operand may still raise
  the inexact exception: the sum of two empty intervals does where GAOL
  checks the rounding direction with an addition. The vector comparisons of
  64-bit ARM (\code{fcmgt}, \code{fcmge}) signal the invalid-operation
  exception on a NaN: the NEON operations tell the empty set, the infinite
  bounds and the zeros from the bits of the bounds, and compare no bound as a
  double before they know it is no NaN.\newinvfive
\item A compiler may compile a quiet comparison into a signaling one, or
  compute a comparison that follows a test of the empty set before it, and
  GAOL is written against the cases seen. GCC turns a choice made on a quiet
  comparison into a conditional move with a signaling comparison on 32-bit
  ARM (\code{vcmpe}, GCC 12 to 14; GCC bug 52258,
  \url{https://gcc.gnu.org/bugzilla/show_bug.cgi?id=52258}, is of this kind)
  and on POWER9 (\code{xscmpgedp}, GCC 13 with \option{-mcpu=power9}): there
  \code{is\_empty()}, the intersection and the constructors of the FPU
  intervals tell NaN bounds with \code{std::isunordered()} first, which stays
  quiet, so that the choice a program makes on \code{is\_empty()}
  (\code{x.is\_empty() ? a : b} for an empty \code{x}), \code{x~\&=~y} for an
  empty \code{y} and \code{floor()}, \code{max()}\dots{} of an empty interval
  raise nothing (the armhf jobs of the continuous integration and a POWER9
  job under qemu check them, built for size too, where GCC calls the
  intersection rather than inlining it). GCC~9.4 at \option{-O3} computed the
  comparisons of the bounds of \code{x~\&=~y} before its test of the empty
  set, in a loop of intersections: they are quiet ones. And GCC vectorizes
  loops of comparisons of bounds into vector comparisons that raise the
  exception for an empty element: GCC~9.4 and~13 where AVX is on
  (\option{-mavx}, or \option{-mfma}, which the three builds give where the
  processor has FMA) and GCC 13 to 15 for 32-bit x86 with SSE2 alone
  (MinGW-w64) make the quiet comparisons signaling predicates
  (\code{vcmpnlepd}, \code{cmpnlepd}\dots), GCC~14 for POWER8 uses
  \code{xvcmpgedp}, and GCC 11 to 15 for 64-bit ARM compute the second and
  third comparisons of the constructor for an empty element too. With GCC and
  the FPU or the NEON intervals, on every processor, the bounds
  \code{is\_empty()} compares, those of the relations that do not call it
  (\code{straddles\_zero()}, \code{set\_contains(d)}\dots) and the doubles
  given to the constructors go through an empty asm statement, which keeps
  GCC from vectorizing such loops: loops of \code{x.set\_contains(d)} and of
  \code{x.straddles\_zero()} took 0.8 and 1.3~ns per element rather than
  0.5~ns with GCC~13 on x86-64, about as long as with GCC~9.4, which did not
  vectorize them and took as long as before. A choice the program makes on a
  relation, \code{x.set\_contains(y) ? a : b}, and a loop of relations other
  than those of \file{tests/rounding\_direction.cpp} are not known to raise
  the exception, but nothing keeps a compiler from compiling them
  so.\newinvfive
\item Nonempty operands raised the invalid-operation exception as well, and
  the cases found no longer do. With the SSE2 intervals, the products
  \code{interval(0) * interval::one\_plus\_infinity()} and
  \code{interval::positive() * interval(0)}, which multiply a zero bound by an
  infinite one, did (the FPU intervals gave the same products without it),
  and so did, through them, \code{pow} of the base \code{interval(1)} and the
  exponent \code{interval::one\_plus\_infinity()}; \code{pow} did in every
  build, through the \code{pow} of CORE-MATH, for an exponent of extreme
  magnitude (\code{pow(interval(1, 2), interval(4.9e-324))},
  \code{pow(interval(1, 2), interval(1e-300))}, an exponent interval with a
  bound near \code{DBL\_MAX} such as \itv{10^{300}}{10^{308}}), which the
  CORE-MATH of GAOL v5 no longer does, and for a power next to the largest
  double (\code{pow(interval(2), interval(1023.9999))},
  \code{pow(interval(0x1.000003f82a1fp+0), 3000000001LL)}), where CORE-MATH
  multiplied its first approximation by $2^{1024}$, an infinity, and computed
  $\infty-\infty$ (changed in the vendored sources, see \file{3rd/README.md});
  and with the SSE2 intervals, \code{div\_rel(K, J, I)} did where it divided
  an infinite bound of \code{K} by an infinite bound of \code{J} in a half of
  the register it did not keep ($K=\itv{2}{+\infty}$ and
  $J=\itv{-\infty}{2}$), and raised the division-by-zero exception where it
  divided a bound by a zero one there ($K=\itv{-2}{-2}$ and $J=\itv{-2}{0}$),
  and so did the division, the relational division and the inverse of a
  divisor with a zero bound (\code{interval(1, 2)/interval(0, 2)},
  \code{interval(1, 2)\%interval(-2, 0)}, \code{inverse(interval(0, 2))},
  \code{pow(interval(0, 2), -1)}). This list is not exhaustive.\newinvfive
\item The overflow and the division-by-zero exceptions are no longer raised
  for a result that has no infinite bound. Some functions computed beyond the
  largest double on the way, or divided by zero, and raised them for a finite
  result, the empty set included: \code{sqrt(interval(1,~DBL\_MAX))},
  \code{nth\_root(interval(DBL\_MAX),~3)} and the roots of order 4 and more,
  and the relational \code{sqrt\_rel()} and \code{nth\_root\_rel()} (the
  product that proves a root exact is $2^{1024}$ for the square root of
  \code{DBL\_MAX}), \code{sin(interval(-1,~DBL\_MAX))} and \code{cos()} (the
  width of the interval), \code{tan(interval(DBL\_MAX))} ($x+\pi/2$),
  \code{div\_rel()} (the quotient by a divisor next to $2^{-1074}$),
  \code{atan\_rel()}, \code{atan2\_rel\_x()} and \code{atan2\_rel\_y()},
  \code{pow\_rel()} and \code{pow\_exponent\_rel()} (the powers beyond the
  largest double that their searches try, and the exponents next to
  $2^{-1074}$), and \code{atanh\_rel()} (atanh at $-1$ and $1$). A quotient,
  a product or a power that would overflow is the infinity, or the largest
  double, that the operation gives, without being computed, where it is only
  an intermediate value: the exceptions still tell an infinite bound of the
  result, as above. This list is not exhaustive either, and the tests check
  these functions and the others of the same kind on operands next to the
  limits of the doubles (\file{tests/rounding\_direction.cpp}).\newinvfive
\end{itemize}

\noindent
An empty operand, and the nonempty operands tested, raising no
invalid-operation exception therefore does not make that exception safe to
leave enabled: every exception stays disabled while GAOL computes.\newinvfive

\medskip
The flags tell nothing of the results: after most operations of GAOL,
\code{fetestexcept(FE\_INEXACT)} is raised whatever the result (not after
\code{-x}, \code{abs()}, \code{\&}, \code{|}, \code{floor()}, \code{max()} or
\code{min()}, which round nothing and check no rounding direction, nor after
\code{+}, \code{-}, \code{/} and \code{sqrt()} with the NEON intervals when
their bounds are exact), and \code{FE\_OVERFLOW} or \code{FE\_DIVBYZERO}
stands for an infinite bound of the interval, or an infinite result of
\code{width()} or \code{rad()}, not for an infinite number the program
computed: \code{midpoint()} and \code{mid()} of
\itv{\text{DBL\_MAX}}{\text{DBL\_MAX}}, which overflowed the sum of the
bounds, raise none. IEEE 1788-2015 \cite{IEEE1788:2015} leaves the flags of
its operations unspecified. A program that reads them for its own
computations clears them (\code{std::feclearexcept(FE\_ALL\_EXCEPT)}) just
before, and reads them before its next operation of GAOL.\newinvfive

\medskip
The initialization of GAOL, which runs before \code{main()}, sets the default
environment (\code{fesetenv(FE\_DFL\_ENV)}), which masks every exception,
unless GAOL is built with \option{GAOL\_PRESERVE\_ROUNDING}: the
initialization then leaves the whole environment as it found it, and an
exception a static object enabled before it stays enabled. A program enables
the exceptions in \code{main()} or later, whichever the build, not in a
static object initialized before GAOL. \code{cleanup()} leaves the exceptions
and the flags as they are, and so does every operation, whichever way GAOL is
built: it leaves the exceptions the program enabled enabled, and the flags it
raised raised, and adds the flags it raises itself. With
\option{GAOL\_PRESERVE\_ROUNDING} and the SSE2 intervals, \code{+}, \code{-},
\code{*}, \code{/}, \code{\%}, \code{div\_rel()}, \code{sqr()},
\code{inverse()}, the integer powers and \code{+=~d}, \code{-=~d},
\code{*=~d}, \code{/=~d}, \code{\%=~d} wrote the SSE control register whole,
every exception masked and every flag cleared: an exception the program had
enabled was masked again by the first of these operations, and the flags it
had raised were lost (\file{tests/rounding\_direction.cpp} checks
it).\newinvfive

%%------------------------------------------------------------------------
\section{Initialization and cleanup}\label{chap:init}
%%------------------------------------------------------------------------
\cindex{initialization}\cindex{cleanup}%

GAOL initializes itself when the program starts, before \code{main()}:
\code{init()} need not be called. Since Release 1.0 of GAOL, there is an
automatic initialization and cleanup, so that GAOL works even when the program
calls neither function. Since GAOL v5, the initialization is made with every
compiler, Visual C++ included: each file including the headers of GAOL holds
a static object whose constructor initializes GAOL the first time, as
\file{<iostream>} initializes the standard streams, so that GAOL is
initialized before the static objects that such a file defines after its
\code{\#include}. Each operation also sets the rounding direction upward
itself (\pxref{sec:rounding-direction}), so that \code{init()} no longer has
to be called again after the calling code changed it.\newinvfive

\begin{Rem}\label{rem:cleanup-rounding}
\code{cleanup()}, on the other hand, is to be called by the program right
after its last use of GAOL, at the latest at the end of \code{main()}, with
every compiler.\newinvfive{} Unless GAOL is built to preserve the rounding
direction, the initialization leaves the rounding direction upward for the
whole program, not only for the operations of GAOL, and \code{cleanup()} sets
back the one the program started with. The automatic cleanup does so too, but
only when the program ends, once \code{main()} has returned and the static
objects are destroyed, or when the library is unloaded: until then, the code
after the last use of GAOL (the rest of \code{main()}, the destructors of
static objects, the functions registered with \code{std::atexit()}) computes
its doubles rounded upward.
\end{Rem}

\begin{Ex}\label{ex:cleanup}
The program below computes an interval, calls \code{cleanup()} right after
its last use of GAOL, and then prints a double of its own, which is computed
to nearest. It calls no \code{init()}.
\begin{example}
#include <cstdio>
#include <iostream>
#include <gaol/gaol.h>
using namespace gaol;

int main()
{
  @rem^// No gaol::init(): GAOL initialized itself before main()~
  interval x(1, 2);
  std::cout << pow(x, 3) + sin(x) << std::endl;

  gaol::cleanup();   @rem^// always, right after the last use of GAOL~
  std::printf("%g\n", 1.0 / 3.0);   @rem^// the program's own doubles, to nearest~
  return 0;
}
\end{example}
\end{Ex}

\medskip
Only the first call of \code{cleanup()} sets the direction back: an operation
of GAOL after it sets the direction upward again, and
\code{restore\_rounding()} sets it back as many times as the program
needs.\newinvfive{} What the initialization allocated is freed by the
automatic cleanup only, after the static objects are destroyed, so that the
expressions of \file{gaol/parser/gaol\_expression.h} still alive when the
program calls \code{cleanup()} remain valid.

\medskip
The direction set back is the one GAOL found when it initialized itself. It
is not the one the program started with when an operation of GAOL came
first: in the initialization of a static object of a file that uses GAOL
without including its headers, through a library of the program for
instance, or, with Clang, of an \code{inline} variable or a static member of
a class template, whose initialization C++ does not order with the rest of
the file.

\medskip
In a build that takes the AVX-512 path (\pxref{sec:avx512}), the
initialization stops the program, with a message on the standard error, when
the processor or the system does not run the AVX-512 instructions, even in a
program that computes nothing.\newinvfive

\medskip
The constants of GAOL ($\pi$, the masks of its SSE2 operations) are
initialized when compiling: a static object of the program may compute
intervals, before the files of GAOL are initialized, which with the static
library of the CMake build come after those of the program, except with
MinGW-w64.\newinvfive

\begin{deffun}{bool}{init}{(int \var{dbg\_lvl} = 0)}
Initializes the variable \code{debug\_level} (\pxref{sec:debugging-facilities})
to the value of \var{dbg\_lvl}.

Unless GAOL was built to preserve the rounding direction, the first call notes
the rounding direction it finds, which \code{cleanup()} sets back, then sets
the floating-point environment to its default one, then the rounding
direction towards $+\infty$. In addition, it sets the number of digits to
display for interval bounds to 16. In a build that takes the AVX-512 path,
the first call stops the program, with a message on the standard error, when
the processor or the system does not run the AVX-512 instructions, even in a
program that computes nothing.\newinvfive

Returns \code{true} if the library was not already initialized and
\code{false} otherwise. After its first call, the function sets the debugging
level to its new value and, unless the rounding direction is preserved, the
rounding direction upward again: the C runtime of Windows resets it when the
program starts, after the library initialized itself.
\end{deffun}

\begin{deffun}{bool}{cleanup}{(void)}
Unless GAOL was built to preserve the rounding direction, sets back the
rounding direction that the first call of \code{init()} found: the one of the
program when GAOL initialized itself, to nearest when the program starts. On
x86 processors, the x87 unit and the SSE instructions each get theirs back.
The rest of the floating-point environment is left as it is, the exception
flags raised since in particular. The rounding direction being one of each
thread, it is set back in the thread calling \code{cleanup()}.\newinvfive

The program calls it right after its last use of GAOL: an operation of GAOL
after it sets the rounding direction upward again. The automatic cleanup
calls it when the program ends, or when the library is unloaded, if the
program did not, and frees what the initialization allocated, after the
static objects of the program are destroyed: the expressions of
\file{gaol/parser/gaol\_expression.h} still alive when the program calls
\code{cleanup()} remain valid.

Returns \code{true} the first time it is called and \code{false} afterwards:
only the first call sets back the rounding direction.
\end{deffun}

\begin{example}
#include <gaol/gaol.h>

int main(void)
{
  gaol::init(1); @rem^// First level of debugging requested;~
                 @rem^// without it, GAOL initializes itself~

  @emph^[Some code using interval arithmetic]~

  gaol::cleanup(); @rem^// Always, right after the last use~
                   @rem^// of GAOL: sets back the rounding~
                   @rem^// direction~
  return 0;
}
\end{example}

\begin{deffun}{void}{restore\_rounding}{(void)}
Sets back the rounding direction that the first call of \code{init()} found,
as \code{cleanup()} does at its first call, but as many times as the program
calls it.\newinvfive{} A program that goes back to GAOL after
\code{cleanup()}, whose operations set the direction upward again, calls it
after each use. Like \code{cleanup()}, it sets back the direction of the
calling thread only, and leaves the rest of the floating-point environment as
it is. With \option{GAOL\_PRESERVE\_ROUNDING} it does nothing, GAOL leaving
the direction of the program alone (\pxref{sec:upward-effects}).
\end{deffun}

%%------------------------------------------------------------------------
\section{The rounding functions of the FPU}\label{chap:fpu}
%%------------------------------------------------------------------------
\cindex{rounding!functions}%

GAOL provides functions to set the rounding direction of the floating-point
unit. The operations of GAOL do not need them: each sets the rounding
direction upward itself (\pxref{sec:rounding-direction}).

\medskip
On x86 processors, with GCC, Clang and Visual C++, they write the control
registers of the x87 and SSE units directly (\code{fnstcw}/\code{fldcw},
\code{stmxcsr}/\code{ldmxcsr}), rather than calling \code{fesetround()},
which cost 50 to 250~ns per call with the C runtime of Windows; with Visual
C++ for x64, the SSE register only, the x87 unit being unused there
(\pxref{sec:x86-control-registers}).\newinvfive{} Elsewhere (ARM, POWER,
s390x, RISC-V), they call \code{fesetround()}, except on 64-bit ARM
processors with the NEON intervals, and on 32-bit ARM processors with GCC
(\pxref{sec:neon}), where they write the rounding direction field of the
control register, FPCR or FPSCR, themselves, and only when the direction
changes.\newinvfive

\begin{deffun}{void}{round\_downward}{(void)}
Sets the rounding direction mode towards $-\infty$.
\end{deffun}

\begin{deffun}{void}{round\_nearest}{(void)}
  Sets the rounding direction mode to the nearest/even.
\end{deffun}

\begin{deffun}{void}{round\_upward}{(void)}
  Sets the rounding direction mode to $+\infty$.
\end{deffun}

%%------------------------------------------------------------------------
\section{Discussion}
%%------------------------------------------------------------------------

This chapter has established how GAOL v5 handles the rounding direction: each
operation that computes a bound checks it with an addition and sets it upward
when it is not, so that the bounds are right whatever direction the calling
code left, where GAOL~4 trusted the user to set it. The price of this choice
falls on the doubles of the program, which are rounded upward from the
initialization of GAOL until \code{cleanup()}: the quotients, the decimal
conversions of the C library, the sums and the error-free transformations all
change, and the program answers by calling \code{cleanup()} right after its
last use of GAOL, by reading and writing its numbers through GAOL, and by
computing to nearest in a guarded block what needs it. The build with
\option{GAOL\_PRESERVE\_ROUNDING} removes the price from the program and puts
it on the operations, 1.2 to 1.6 times slower with the AVX-512 path and
several times without it; the control registers are then written only where
they have to be, and the AVX-512 instructions, whose rounding direction is in
the instruction, are taken where the machine has them.

\medskip
The same check defends the bounds against the modes that flush the subnormal
numbers to zero, which a program linked with \option{-Ofast} gets without
asking, and which GAOL clears on x86 and on ARM, while \option{-mno-daz-ftz}
keeps them out of the link where the compiler accepts it. The floating-point
exceptions, on the other hand, are left to the program, which disables them
while GAOL computes: the operations raise the inexact flag almost always, the
overflow and division-by-zero flags for an infinite bound, and no longer the
invalid-operation one for an empty operand in the cases found, but nothing
keeps a compiler from compiling a quiet comparison into a signaling one.
Several points stay open in the sources: the other processors (POWER, s390x,
RISC-V, MIPS) have no defence against their flush modes, and nothing was
tested with them set; the FPU intervals of POWER, s390x, RISC-V and
MinGW-w64 for arm64 are not measured; the slower line of 13 operations with
the direction preserved on Visual C++ for x86 has no explanation yet; and the
processors that do not halve their AVX-512 instructions are expected, not
measured, to take less time on that path.

\medskip
With the floating-point environment in place, the next chapters turn to the
intervals themselves: Chapter~\ref{chap:creation} describes how an interval
is created, from numbers or from text, and assigned.
